% frontmatter/zdania/mapa-del-curso.tex -- el año entero en una tabla,
% como frontmatter/mapa-del-curso.tex, en polaco (la última columna sin
% justificar: con palabras polacas, largas, se llenaba de huecos); y,
% debajo, las letras de "Pisz po śladzie", que en polaco son 32.
\section*{Mapa roku}

\begin{center}
\renewcommand{\arraystretch}{1.4}
\begin{tabularx}{\linewidth}{@{}c c c c >{\raggedright\arraybackslash}X@{}}
\toprule
\textbf{Część} & \textbf{Dni} & \textbf{Pora roku} & \textbf{Zdań dziennie} & \textbf{Co się dzieje} \\
\midrule
1 & 1--65   & jesień & 1 & Poznajemy Lucíę i jej rodzinę: Daniego,
Toby'ego, babcię Rosę. Część kończy się Bożym Narodzeniem. \\
2 & 66--130 & zima & 2 & Trzech Króli, śnieg, karnawał -- i już czuć
wiosnę. \\
3 & 131--195 & wiosna & 3 & Małe historie w trzech zdaniach, Dzień
Książki i koniec roku szkolnego. \\
4 & 196--260 & lato & 4 & Wakacje, dłuższe historie, trochę dialogów i
wielki finał. \\
\bottomrule
\end{tabularx}
\end{center}

\vspace{6mm}
Każda część ma 13 tygodni po 5 dni -- także w czasie przerwy
świątecznej i wakacji: zeszyt to wciąż jedna strona na każdy dzień od
poniedziałku do piątku, w Boże Narodzenie i latem, nie tylko w szkole.
Liczba zdań rośnie raz, na początku każdej części -- w tej samej części
wszystkie strony mają tyle samo zdań.

Mniej więcej co dwa tygodnie, \emph{\lblTraza}: litery polskiego
alfabetu po kolei, osiem w każdej części -- od \emph{a} jesienią do
\emph{ż} latem.
\newpage
